\section*{Bab \ref{ch:g5:solids} --- Kubus dan Balok}

\begin{solution}{exo:g5:solids:1}
Misalnya: sebuah dadu (\omterm{def:g5:solids:def}{kubus}); kotak sereal (\omterm{def:g5:solids:def}{balok}); sebuah kaleng
(\omterm{def:g5:solids:def}{tabung}); bola sepak (bola); atap tertentu atau monumen di Mesir
(\omterm{def:g5:solids:def}{limas}).
\end{solution}

\begin{solution}{exo:g5:solids:2}
Sebuah \omterm{def:g5:solids:def}{balok} punya $6$ sisi, $12$ \omterm{def:g5:solids:def}{rusuk}, $8$ \omterm{def:g5:solids:def}{titik sudut} --- dan
$6 + 8 = 12 + 2$. \checkmark
\end{solution}

\begin{solution}{exo:g5:solids:3}
Keenam sisinya terdiri atas tiga pasang sisi berhadapan yang
identik (atas--bawah, depan--belakang, kiri--kanan). Pada \omterm{def:g5:solids:def}{kubus},
keenam sisinya adalah persegi yang identik.
\end{solution}

\begin{solution}{exo:g5:solids:4}
\omterm{def:g5:solids:net}{Jaring-jaring} berbentuk salib: empat persegi $2$ cm dalam satu
baris, satu di atas dan satu di bawah persegi yang kedua. Melipat
cincin berisi empat itu membuat sisi tegaknya; dua persegi tambahan
menutup atas dan bawahnya.
\end{solution}

\begin{solution}{exo:g5:solids:5}
Pada \omterm{def:g5:solids:net}{jaring-jaring} salib (empat persegi dalam satu baris, satu di
atas, satu di bawah persegi kedua), pasangan yang berhadapan adalah:
persegi $1$ dan $3$ pada baris itu; persegi $2$ dan $4$ pada baris
itu; serta persegi atas dan bawah. Satu penamaan yang benar: baris
$= 1, 2, 6, 5$; atas $= 3$; bawah $= 4$ --- lalu $1{+}6$, $2{+}5$,
$3{+}4$ semuanya menjadi $7$.
\end{solution}

\begin{solution}{exo:g5:solids:6}
\omterm{def:g5:solids:def}{Balok}: $5 \times 2 = 10$ per lapisan, $3$ lapisan: $30$ \omterm{def:g5:solids:def}{kubus}.
\omterm{def:g5:solids:def}{Kubus}: $3 \times 3 \times 3 = 27$. Tumpukan L:
$3 \times 2 \times 1 = 6$ ditambah $1 \times 2 \times 1 = 2$: $8$ \omterm{def:g5:solids:def}{kubus}.
\end{solution}

\begin{solution}{exo:g5:solids:7}
Lapisan berisi $12$: misalnya $4 \times 3$ atau $6 \times 2$ (atau
$12 \times 1$). Dengan $3$ lapisan, \omterm{def:g5:solids:volume}{volumenya}
$12 \times 3 = 36$ \omterm{def:g5:solids:def}{kubus} dalam setiap keadaan.
\end{solution}

\begin{solution}{exo:g5:solids:8}
$27 = 3 \times 3 \times 3$: \omterm{def:g5:solids:def}{kubus} dengan sisi $3$.
$64 = 4 \times 4 \times 4$: \omterm{def:g5:solids:def}{kubus} dengan sisi $4$.
\end{solution}

\begin{solution}{exo:g5:solids:9}
Kotaknya memuat $10 \times 5 \times 4 = 200$ butir gula batu. Dengan
$6$ butir sehari: $200 = 6 \times 33 + 2$, jadi kotaknya bertahan
$33$ hari penuh (dan $2$ butir tersisa untuk pagi hari ke-$34$).
\end{solution}

\begin{solution}{exo:g5:solids:10}
Terkena cat pada $3$ sisi: $8$ \omterm{def:g5:solids:def}{kubus} di pojoknya. Pada $2$ sisi:
bagian tengah $12$ \omterm{def:g5:solids:def}{rusuknya}, yaitu $12$ \omterm{def:g5:solids:def}{kubus}. Pada $1$ sisi: pusat
dari $6$ sisinya, yaitu $6$ \omterm{def:g5:solids:def}{kubus}. Tidak sama sekali: satu \omterm{def:g5:solids:def}{kubus} pusat
yang tersembunyi, yaitu $1$. Periksa: $8 + 12 + 6 + 1 = 27$. \checkmark

\end{solution}
